% Input from reports/phase2/ with amsmath and
% cite loaded. New keys: sections/theory.bib; other keys are listed there.
% Model settings quoted here come from configs/{lstm,tcn,transformer,hybrid_gated,dl_base}.yaml
% and src/rul/models/; numeric illustrations were computed for this section
% (receptive fields from rul.models.tcn.receptive_field; NASA-score shift by
% simulation with Gaussian errors, stated as such).

\section{Theory of the Phase 2 Models}
\label{sec:theory}

All four sequence models read the same input: a window $X\in\mathrm{R}^{L\times F}$
of the last $L=30$ cycles of $F$ scaled sensors, labelled with the capped RUL
of its last cycle. They differ only in how information moves across the $L$
cycles, and that difference is what this section explains.

\subsection{LSTM: Vanishing Gradients and Gates}
A plain recurrent network updates $h_t=\tanh(Wh_{t-1}+Ux_t+b)$. Training it
by backpropagation through time multiplies one Jacobian per step,
\begin{equation}
\frac{\partial h_L}{\partial h_t}=\prod_{s=t+1}^{L}\operatorname{diag}\!\bigl(\tanh'(\cdot)\bigr)\,W ,
\end{equation}
so the gradient from cycle $t$ shrinks (or grows) roughly like a power of
$\lVert W\rVert$~\cite{Bengio1994learning,Pascanu2013difficulty}. With a
factor of 0.9 per step, the first cycle of a 30-cycle window gets
$0.9^{29}\approx0.05$ of the signal of the last one. Slow wear is exactly a
long-range pattern, so a plain RNN learns it poorly.

The LSTM~\cite{Hochreiter1997long} adds a cell state $c_t$ and three gates,
each a sigmoid in $[0,1]$ computed from $x_t$ and $h_{t-1}$:
\begin{align}
c_t &= f_t\odot c_{t-1} + i_t\odot\tilde c_t, & h_t &= o_t\odot\tanh(c_t),
\end{align}
with forget gate $f_t$, input gate $i_t$, output gate $o_t$ and candidate
$\tilde c_t=\tanh(W_cx_t+U_ch_{t-1}+b_c)$. Along the cell path,
$\partial c_t/\partial c_{t-1}=\operatorname{diag}(f_t)$ (plus indirect
terms). The gradient is therefore scaled by a forget gate the network learns,
not by a fixed weight matrix: when $f_t\approx1$ the memory, and its gradient,
pass through almost unchanged~\cite{Gers2000learning}. The input gate decides
when a new reading is worth writing, which suits C-MAPSS, where most cycles
are noise around a slow trend. Our LSTM has one layer of 64 units (about
21,000 weights in the recurrent layer for $F=15$) and uses $h_L$ as the
window's embedding. The gates do not stop gradients from \emph{exploding},
which is why the trainer clips the gradient norm at 1.

\subsection{Temporal CNN: Dilated Causal Convolutions}
A temporal convolutional network (TCN)~\cite{Bai2018empirical} applies 1-D
convolutions along time. \emph{Causal} padding (zeros on the left only)
means the output at cycle $t$ depends only on cycles $\leq t$, so no future
reading leaks in. \emph{Dilation} $d$ spaces the $k$ filter taps $d$ cycles
apart: $y_t=\sum_{j=0}^{k-1}w_j\,x_{t-jd}$. Our residual blocks hold two such
convolutions with dilation $2^{\ell}$ at level $\ell=0,\dots,N-1$. Each
convolution extends the view by $(k-1)d$ cycles, so the last output sees
\begin{equation}
R = 1 + 2(k-1)\sum_{\ell=0}^{N-1}2^{\ell} = 1+2(k-1)(2^{N}-1)
\end{equation}
cycles. The field doubles with each level while the weight count grows only
linearly. With $k=3$, two levels see 13 cycles, three see 29 and four see 61.
Three levels would therefore ignore the oldest cycle of our 30-cycle window,
which is why the configuration uses four. Unlike the LSTM, the TCN's
``memory'' is a hard, known number, and every cycle inside it is processed in
parallel.

\subsection{Transformer: Attention, Positions and Size}
Self-attention~\cite{Vaswani2017attention} projects each cycle's vector to a
query $q_t$, key $k_t$ and value $v_t$, and replaces it with a weighted
average of all values,
\begin{equation}
\operatorname{Attn}(Q,K,V)=\operatorname{softmax}\!\Bigl(\frac{QK^\top}{\sqrt{d_k}}\Bigr)V ,
\end{equation}
so any two cycles interact in one step, whatever their distance. The
$\sqrt{d_k}$ keeps the dot products from saturating the softmax. Several
heads run in parallel and can look for different patterns. Attention is
order-blind: shuffling the cycles shuffles the output in the same way. A
sinusoidal code is therefore added to each cycle's vector,
$\mathrm{PE}_{t,2i}=\sin(t/10000^{2i/d})$ and
$\mathrm{PE}_{t,2i+1}=\cos(t/10000^{2i/d})$, which gives every position a
unique pattern that varies smoothly in $t$.

\emph{Why size matters here.} Attention assumes very little about the data
(no recurrence, no locality), so it must learn that structure from
examples. FD001 offers about 20,000 windows, but they come from only 100
engines, and neighbouring windows overlap in 29 of their 30 cycles. The
effective sample size is much closer to 100 trajectories than to 20,000
points. A large Transformer would memorize the training engines, so ours
keeps $d=32$, 4 heads, one pre-norm layer and a feed-forward width of 64,
roughly 9,000 encoder weights, which is fewer than the LSTM.

\subsection{Gated Fusion}
The gated hybrid embeds the window with the LSTM ($h_{\mathrm{seq}}$) and the
engineered features at the window's last cycle with a small MLP
($h_{\mathrm{feat}}$), both of size $d$, and mixes them per
dimension~\cite{Arevalo2017gated}:
\begin{align}
g&=\sigma\bigl(W[h_{\mathrm{seq}};h_{\mathrm{feat}}]+b\bigr),\\
z&=g\odot h_{\mathrm{seq}}+(1-g)\odot h_{\mathrm{feat}},\qquad \hat y=\mathrm{head}(z).
\end{align}
Unlike plain concatenation~\cite{Chen2021machine,Jin2022bilstm}, the mix can
change from one window to the next. For example, the network may lean on the
rolling features early in life and on the sequence near failure.

\emph{What the gate can tell us:} for each window, how much each embedding
dimension is drawn from each branch, so its average by life stage is a
useful description. \emph{What it cannot tell us:}
(i)~$g$ is a mixing weight, not a contribution. A dimension with $g=0.9$ but
$h_{\mathrm{seq}}\approx0$ adds almost nothing, and the head can amplify or
ignore any dimension.
(ii)~The two branches are trained together and can co-adapt, so a feature
branch that looks unused may still shape what the sequence branch learns.
(iii)~Averaging $g$ over dimensions hides structure.
(iv)~Like attention weights~\cite{Jain2019attention}, gate values are not
causal explanations. The causal check is the ablation: retrain with one
branch removed (\texttt{fusion: sequence} or \texttt{features}) and compare
the errors.

\subsection{Regularization and Bias--Variance}
For squared error, the expected test error at an input splits
into~\cite{Geman1992neural}
\begin{equation}
\mathrm{E}\bigl[(y-\hat f(x))^2\bigr]
=\underbrace{\bigl(f(x)-\mathrm{E}\hat f(x)\bigr)^2}_{\text{bias}^2}
+\underbrace{\operatorname{Var}\hat f(x)}_{\text{variance}}
+\sigma^2 ,
\end{equation}
where the expectation is over training sets. On C-MAPSS, $\sigma^2$ is large
early in life, because a healthy engine's sensors do not reveal its RUL. The
cap removes most of that irreducible part from the target. Deep sequence
models have low bias but high variance on 100 engines: refitting on another
fold, or with another seed, moves their predictions noticeably. Each
regularizer trades a little bias for less variance:
dropout~\cite{Srivastava2014dropout} trains an implicit ensemble of thinned
networks; weight decay shrinks weights towards zero; input
noise~\cite{Bishop1995noise} acts like a smoothness penalty; and early
stopping~\cite{Prechelt1998early} limits how far the optimizer moves from
small initial weights, judged on engines held out of training. In our study
(Section~\ref{sec:dl_dataset}) none of these beat the baseline, which already
had dropout 0.3 and early stopping. The fold standard deviation measures
variance across data splits, but not across seeds, so repeated seeds are
still needed to separate the two.

\subsection{Loss Choice and the NASA Score}
Training minimizes the mean squared error, whose minimizer is the
conditional mean $\mathrm{E}[y\mid X]$: it treats a prediction 20 cycles too
late the same as one 20 cycles too early. The PHM08
score~\cite{Saxena2008damage} does not. With error $e=\hat y-y$ it adds
$e^{-e/13}-1$ for early and $e^{e/10}-1$ for late predictions, and it is a
sum over engines, not a mean. At $|e|=20$ a late prediction costs 6.4 against
3.7 for an early one. At $|e|=40$ it costs 53.6 against 20.7. The score is
therefore dominated by a few large late errors.

Because of this asymmetry, the best prediction under the NASA score lies
below the conditional mean. As an illustration, if errors were Gaussian
with standard deviation 16 cycles, shifting every prediction 4.5 cycles
earlier would lower the expected score per engine from 4.4 to 3.9 (about
11\%), while RMSE would rise only from 16.0 to 16.6. That is why we report
RMSE, MAE and the NASA score together. A model selected on RMSE is not
automatically the best on the NASA score. An asymmetric training loss, or a
post-hoc early shift chosen on validation folds, could target the score
directly, at some cost in RMSE.
